% Options for packages loaded elsewhere
\PassOptionsToPackage{unicode}{hyperref}
\PassOptionsToPackage{hyphens}{url}
\documentclass[
]{article}
\usepackage{xcolor}
\usepackage[margin=1in]{geometry}
\usepackage{amsmath,amssymb}
\setcounter{secnumdepth}{-\maxdimen} % remove section numbering
\usepackage{iftex}
\ifPDFTeX
  \usepackage[T1]{fontenc}
  \usepackage[utf8]{inputenc}
  \usepackage{textcomp} % provide euro and other symbols
\else % if luatex or xetex
  \usepackage{unicode-math} % this also loads fontspec
  \defaultfontfeatures{Scale=MatchLowercase}
  \defaultfontfeatures[\rmfamily]{Ligatures=TeX,Scale=1}
\fi
\usepackage{lmodern}
\ifPDFTeX\else
  % xetex/luatex font selection
\fi
% Use upquote if available, for straight quotes in verbatim environments
\IfFileExists{upquote.sty}{\usepackage{upquote}}{}
\IfFileExists{microtype.sty}{% use microtype if available
  \usepackage[]{microtype}
  \UseMicrotypeSet[protrusion]{basicmath} % disable protrusion for tt fonts
}{}
\makeatletter
\@ifundefined{KOMAClassName}{% if non-KOMA class
  \IfFileExists{parskip.sty}{%
    \usepackage{parskip}
  }{% else
    \setlength{\parindent}{0pt}
    \setlength{\parskip}{6pt plus 2pt minus 1pt}}
}{% if KOMA class
  \KOMAoptions{parskip=half}}
\makeatother
\usepackage{graphicx}
\makeatletter
\newsavebox\pandoc@box
\newcommand*\pandocbounded[1]{% scales image to fit in text height/width
  \sbox\pandoc@box{#1}%
  \Gscale@div\@tempa{\textheight}{\dimexpr\ht\pandoc@box+\dp\pandoc@box\relax}%
  \Gscale@div\@tempb{\linewidth}{\wd\pandoc@box}%
  \ifdim\@tempb\p@<\@tempa\p@\let\@tempa\@tempb\fi% select the smaller of both
  \ifdim\@tempa\p@<\p@\scalebox{\@tempa}{\usebox\pandoc@box}%
  \else\usebox{\pandoc@box}%
  \fi%
}
% Set default figure placement to htbp
\def\fps@figure{htbp}
\makeatother
\setlength{\emergencystretch}{3em} % prevent overfull lines
\providecommand{\tightlist}{%
  \setlength{\itemsep}{0pt}\setlength{\parskip}{0pt}}
\usepackage{bookmark}
\IfFileExists{xurl.sty}{\usepackage{xurl}}{} % add URL line breaks if available
\urlstyle{same}
\hypersetup{
  pdftitle={Experiments Problem Set 1},
  pdfauthor={Kristina Mensik},
  hidelinks,
  pdfcreator={LaTeX via pandoc}}

\title{Experiments Problem Set 1}
\author{Kristina Mensik}
\date{2026-08-31}

\begin{document}
\maketitle

1a. An experiment is a procedure for estimating a causal quantity of
interest, where treatment is randomized: probability of assignment to
treatment is a known quantity between but not equal to 0 and 1, and can
be replicated, and is equally likely across outcome-relevant observed
and unobserved factors. The difference between an experiment and
observational study is that in the prior, assignment to treatment is
randomized.

Some though argue that to be an experiment, the entire system under
study (including assignment to treatment and manipulations) is under the
researcher's control.

1b. Unobserved heterogeneity is the unmeasured, unknown, potentially
effectively infinite attributes that may distort inference. It is a
problem for that reason (some other set of factors could explain a
correllational relationship, even between variables that are comprable
on observables), but especially a problem because there's no resolution
procedure. In other words, there is no way of knowing what you don't
know.

\begin{enumerate}
\def\labelenumi{\arabic{enumi}.}
\setcounter{enumi}{1}
\item
  This is none of the above -- it is more like an observational study --
  because just because the pool of villages is randomly sampled, does
  not mean that assignment to treatment for the population under study
  is random.
\item
  This study seems to adequately fulfill the criteria for a field
  experiment. First, demand for household water connections is measured
  by the percent of uptake on offers for zero-interest credit household
  water connections -- the treatment resembles the intervention of
  interest. Second, those who would encounter the intervention are
  households in urban Morocco who might want household water connections
  (those who do not have them already), and as long as ``homeowners
  without a private connection to the city's water grid'' do not have
  another form of household water connections, the participants indeed
  resemble the actors of interest. Last, the context of receiving
  treatment also seems to match, again so long as the subset it random,
  and the outcomes do too if waterbourne illness, leisure time and
  social activities are thought of as good indicators of welfare (I
  think so).
\item
  It is reasonable to believe that parachutes are effective even in the
  abasence of randomized experiments that establish their efficacy
  because there is very little ambiguity about the potential outcome for
  the treated in a world where they are not assigned to treatment, nor
  about what the potential outcome would be if there were in fact a
  procedure assigning people to be ``untreated.''
\end{enumerate}

\end{document}
